\documentclass[a4paper,11pt]{article}
% Подключаем все нужные пакеты
%%% Локализация и шрифты
\usepackage[utf8]{inputenc}          % Кодировка UTF-8
\usepackage[T2A]{fontenc}            % Русская кодировка
\usepackage[english,russian]{babel} % Локализация
\usepackage{cmap}                    % Поддержка поиска в PDF
\usepackage{tikz}
\usepackage{xstring}
\usepackage[left=20mm, top=10mm, right=20mm, bottom=10mm, nohead, nofoot]{geometry}
\renewcommand{\baselinestretch}{1.0}% Межстрочный интервал
\usepackage{ragged2e}               % Выравнивание текста
\usepackage{microtype}              % Улучшение качества вывода текста
\usepackage{multicol}
\usepackage{tcolorbox}
\tcbuselibrary{listingsutf8, breakable, skins}
\justifying
\sloppy
\tolerance=500
\hyphenpenalty=10000
\emergencystretch=3em

\usepackage[normalem]{ulem}
\pagestyle{empty} \textwidth=19.0cm \oddsidemargin=-1.3cm
\textheight=26cm \topmargin=-3.0cm
\usepackage{tcolorbox}
\tcbuselibrary{listingsutf8}
\usepackage{pgfplots}
\renewcommand{\rmdefault}{cmss}
\usepgfplotslibrary{patchplots}
\pgfplotsset{compat=1.18}
\usetikzlibrary{fillbetween}
% Создаём новый счётчик задач
\newcounter{mainth}
\newcounter{subth}[mainth]

\newcommand{\thsub}{\arabic{mainth}.\arabic{subth}}

\newcounter{mainpr}
\newcounter{subpr}[mainpr]

\newcommand{\propsub}{\arabic{mainpr}.\arabic{subpr}}

\newcounter{mainex}
\newcounter{subex}[mainex]

\newcommand{\exsub}{\arabic{mainex}.\arabic{subex}}

\renewcommand{\familydefault}{\ttdefault}

% Определяем стили блоков
\usepackage{xcolor}

% --- Однострочное создание блоков \problem, \solution, \answer ---
% (работает без \begin ... \end)
\definecolor{myrose}{RGB}{255,182,193} % Розовый
\definecolor{mybrown}{RGB}{210,180,140} % Светло-коричневый
\definecolor{mycyan}{RGB}{135,206,235}  % Голубой
\definecolor{myorange}{RGB}{255,165,0}   % Оранжевый
\definecolor{mypurple}{RGB}{216,191,216} % Светло-фиолетовый
% Требуется пакет environ:
\newcommand{\thletter}{}

\newcommand{\thtitle}{
  \IfStrEq{\thletter}{}%  % Если taskletter пустая
    {\ttfamily}%   % Выводим только номер
    {\ttfamily (\thletter)}% Иначе номер + буква
}

% --- Определяем сокращённые версии окружений ---

\newtcolorbox{defBox}{
  breakable,
  enhanced,
  colback=violet!10!white,
  colframe=mypurple,
  fonttitle=\ttfamily,
  title={Определение}
}

\newcommand{\defin}[1]{%
  \begin{defBox}%
    #1%
  \end{defBox}%
}

\newtcolorbox{theoremBox}{
  breakable,
  enhanced,
  colback=red!10!white,
  colframe=myrose,
  fonttitle=\ttfamily,
  title={Теорема~\thsub \thtitle}
}

\newcommand{\theorem}[1]{%
  \refstepcounter{subth}%
  \begin{theoremBox}%
    #1%
  \end{theoremBox}%
}

\newtcolorbox{proposBox}{
  breakable,
  enhanced,
  colback=brown!10!white,
  colframe=mybrown,
  fonttitle=\ttfamily,
  title={Утверждение~\propsub}
}

\newcommand{\propos}[1]{%
  \refstepcounter{subpr}%
  \begin{proposBox}%
    #1%
  \end{proposBox}%
}

\newtcolorbox{proofBox}{
  breakable,
  enhanced,
  colback=cyan!10!white,
  colframe=mycyan,
  fonttitle=\ttfamily,
  title={Доказательство}
}

\newcommand{\prof}[1]{
  \begin{proofBox}
    #1
  \end{proofBox}
}

\newtcolorbox{exampleBox}{
  breakable,
  enhanced,
  colback=orange!10!white,
  colframe=myorange,
  fonttitle=\ttfamily,
  title={Пример~\exsub}
}

\newcommand{\example}[1]{%
  \refstepcounter{subex}%
  \begin{exampleBox}%
    #1%
  \end{exampleBox}%
}

% Команда для векторов
\newcommand{\vect}[1]{\mathbf{#1}}
\usepackage[upint]{stix2}

%%% Колонтитулы
\usepackage{fancyhdr}
\pagestyle{fancy}
\renewcommand{\headrulewidth}{0.3mm} % Толщина линии в колонтитулах
% ПРОВЕРЬТЕ ИМЯ ФАЙЛА! Должно быть 'profimatika' или 'profimatica'

\rhead{Теория вероятностей и математическая статистика}
\lhead{Благий Александр}
\headsep=25pt
\footskip=15mm
\textheight=700pt
\headheight=35pt % Возможно, понадобится увеличить до ~30pt. Проверьте log-файл.
\parindent=0mm
%%% Математика
\usepackage{amsmath, amsfonts, amssymb, amsthm, mathtools} % Базовые математические пакеты
\usepackage{icomma}   
\usepackage{euscript}               % Шрифт Евклид
\usepackage{mathrsfs}               % Красивый матшрифт
\usepackage{cancel}                 % Зачёркивание в формулах

% Замены часто употребляемых шрифтов
\newcommand{\mb}[1]{\mathbb{#1}}
\newcommand{\mc}[1]{\mathcal{#1}}
\newcommand{\ms}[1]{\mathscr{#1}}
\begin{document}
\setcounter{maindf}{3}
\setcounter{subdf}{0}

\setcounter{mainth}{3}
\setcounter{subth}{0}

\setcounter{mainpr}{3}
\setcounter{subpr}{0}

\setcounter{mainex}{3}
\setcounter{subex}{0}
\begin{center}
    \textbf{\Huge{Теория вероятностей и математическая статистика}}
    \vspace{5cm}

    Конспекты лекций Коссовой Е.В.

    Факультет экономических наук, 2025/2026 гг.

    \vspace{5cm}

    Особую благодарность хочу выразить студентке ЭАДа Доможаковой Валентине, на основе именно её конспектов составлен данный файл.

    \vspace{10cm}
    \textit{Примечание: данные конспекты начинаются примерно с октября 2025 года, за материалом до этого времени обращайтесь к конспектам Д. А. Борзых}

\end{center}
\newpage
\begin{center}
    \huge{Вероятностная мера}
\end{center}
$\Omega$, $\sigma$-алгебра $\mc{F}$ подмножеств множества $\Omega$

Любое подмножество $ A \in \mc{F}$ называется \textbf{событием}. Элементы $\Omega$ называются \textbf{элементарными исходами}.
\defin{
    Пусть задано измеримое пространство $(\Omega,\, \mc{F})$. Функция $P: \mc{F} \to [0, 1]$ называется
    
    \underline{вероятностной мерой}, если выполняются условия:
    \begin{enumerate}
        \item $P(\Omega) = 1$ (условие нормировки)
        \item Аксиома счётной аддитивности:
        Если $A_1, A_2, \ldots \in \mc{F}\, -$ попарно непересекющиеся множества,
        
        то
        \[
            P\left( \bigcup_{i=1}^{\infty} A_i \right) = \sum_{i=1}^{\infty} P(A_i).
        \]
    \end{enumerate}
}

\propos{
    \[P(\varnothing) = 0\]
}
\prof{
    Рассмотрим $A_1,\,A_2, \dots A_i = \varnothing$ и пусть $P(A_i) = p > 0$. Тогда
    \[1 \geq P\left(\bigcup_{i = 1}^\infty A_i\right) = \sum_{i = 1}^{\infty} P(A_i) = p + p + \cdots\]
    $\implies$ если $p > 0$, то ряд расходится $\implies p = 0$
}
\propos{
    Для любых попарно непересекющихся множеств $A_i \in \mc{F} \ (1 \leq n)$
    \[P\left(\bigsqcup_{i = 1}^n A_i\right) = \sum_{i = 1}^{n} A_i\]
}
\prof{
    Дополним набор множеств из условия счётным набором $A_{n+1}, A_{n+2}, \dots$, каждое из которых пусто. Тогда по утверждению 3.1 $P(A_i) = 0,\ i > n$
    \[\implies P\left(\bigcup_{i = 1}^\infty A_i\right) = \sum_{i = 1}^{n} P(A_i) + \sum_{i = n+1}^{\infty} P(A_i) = \sum_{i = 1}^{n} P(A_i) + 0 = \sum_{i = 1}^{n} P(A_i)\]
}
\newpage
\propos{
    \[\forall\, A_1,A_2 \in \mc{F},\ A_1\subseteq A_2 \implies P(A_2 \setminus A_1) = P(A_2) - P(A_1)\]
}
\prof{
    \[A_2 = A_1 \cup (A_2 \setminus A_1)\]
    \[\implies P(A_2) = P(A_1) + P(A_2 \setminus A_1)\]
    \[\implies P(A_2 \setminus A_1) = P(A_2) - P(A_1)\]
}
\propos{
    \[\forall\, A_1,A_2 \in \mc{F},\ A_1\subseteq A_2 \implies P(A_1) \leq P(A_2)\]
    \textit{очев}
}
\propos{
    \[\forall\, A \in \mc{F} \implies P(A^c) = 1 - P(A)\]
}
\prof{
    \[\Omega = A \cup A^c\]
    \[\implies P(\Omega) = P(A) + P(A^c)\]
    \[\implies 1 = P(A) + P(A^c)\]
    \[\implies P(A^c) = 1 - P(A)\]
}
\propos{
    \[\forall\, A_1,A_2 \in \mc{F} \implies P(A_1 \cup A_2) = P(A_1) + P(A_2) - P(A_1 \cap A_2)\]
}
\prof{
    \[A_1 \cup A_2 = (A_1 \setminus A_2) \cup (A_1 \cap A_2) \cup (A_2 \setminus A_1)\]
    \[\implies P(A_1 \cup A_2) = \underbrace{P(A_1 \setminus A_2) + P(A_1 \cap A_2)}_{P(A_1)} + \underbrace{P(A_2 \setminus A_1) + P(A_1 \cap A_2)}_{P(A_2)} - P(A_1 \cap A_2) =\]
    \[= P(A_1) + P(A_2) - P(A_1 \cap A_2)\]
}
\newpage
\propos{
    \[\forall A_1,\dots,A_n \in \mc{F} \implies P\left(\bigcup_{i = 1}^n A_i\right) \leq \sum_{i = 1}^{n} P(A_i)\]
}
\prof{
    Введём новый набор множеств:
    \[B_1 = A_1\]
    \[B_2 = A_2 \setminus A_1\]
    \[\vdots\]
    \[B_n = A_n \setminus \bigcup_{i = 1}^{n-1} A_i\]
    Тогда $B_i \subseteq A_i\, \forall i$, $B_i \cap B_j = \varnothing\ \forall i \neq j$ и
    \[\bigcup_{i = 1}^n B_i = \bigcup_{i = 1}^n A_i\]
    Поэтому
    \[P\left(\bigcup_{i = 1}^n B_i\right) = P\left(\bigcup_{i = 1}^n A_i\right) = \sum_{i = 1}^{n} P(B_i) \leq \sum_{i = 1}^{n} P(A_i)\]
    }
\propos{
    \[\forall A_1,\dots \in \mc{F} \implies P\left(\bigcup_{i = 1}^\infty A_i\right) \leq \sum_{i = 1}^{\infty} P(A_i)\]
    \textit{Доказательство аналогично предыдущему утверждению.}
}
\newpage
\renewcommand{\thletter}{о непрерывности вероятностной меры}
\theorem{
    \begin{enumerate}
        \item $\forall\, A_1,\dots, A_n, \dots \in \mc{F},\ A_1 \subseteq A_1 \subseteq \cdots \subseteq A_n \subseteq \cdots$
        \[\lim_{n \to \infty} P(A_n) = P\left(\bigcup_{i = 1}^\infty A_i\right)\]
        \item $\forall\, A_1,\dots, A_n, \dots \in \mc{F},\ A_1 \supseteq A_2 \supseteq \cdots \supseteq A_n \supseteq \cdots$
        \[\lim_{n \to \infty} P(A_n) = P\left(\bigcap_{i = 1}^\infty A_i\right)\]
    \end{enumerate}
}
\prof{
    1. Введём новый набор множеств:
    \[B_1 = A_1\]
    \[B_2 = A_2 \setminus A_1\]
    \[\vdots\]
    \[B_i = A_i \setminus A_{i-1}\ \forall i \geq 2\]
    $B_1,\dots,B_n,\dots$ попарно не пересекаются
    \[\bigcup_{i=1}^\infty B_i = \bigcup_{i=1}^\infty A_i,\ P(B_i) = P(A_i) - P(A_{i-1})\]
    \[\implies P\left(\bigcup_{i=1}^\infty A_i\right) = P\left(\bigcup_{i=1}^\infty B_i\right) = \sum_{i=1}^{\infty}P(B_i) = \lim_{n\to\infty}\sum_{i=1}^{n}P(B_i) =\]
    \[= \lim_{n\to\infty} (\cancel{P(A_1)} + \cancel{P(A_2)} - \cancel{P(A_1)} + \cancel{P(A_3)} - \cancel{P(A_2)} + \cdots + P(A_n) - \cancel{P(A_{n-1})}) = \lim_{n\to\infty} P(A_n)\]
    
    2. $A_1^c \subseteq A_2^c \subseteq \cdots \subseteq A_n^c \subseteq \cdots$

    $\implies \displaystyle\lim_{n\to\infty} P\left(A_n^c\right) = P\left(\bigcup_{i=1}^\infty A_i^c\right)\ -$ следует из пункта 1.
    \[\lim_{n\to\infty}P(A_n^c) = \lim_{n\to\infty}(1 - P(A_n)) = 1 - \lim_{n \to \infty} P(A_n)\]
    \[P\left(\bigcup_{i=1}^\infty A_i^c\right) \overset{\text{ф-ла де Моргана}}{=} P\left(\left(\bigcap_{i=1}^\infty P(A_i)\right)^c\right) = 1 - P\left(\bigcap_{i=1}^\infty A_i\right)\]
    Итого
    \[1 - \lim_{n\to \infty}P(A_n) = 1 - P\left(\bigcap_{i=1}^\infty A_i\right)\]
    \[\lim_{n \to \infty} P(A_n) = P\left(\bigcap_{i=1}^\infty A_i\right)\]
}
\newpage
\setcounter{maindf}{4}
\setcounter{subdf}{0}

\setcounter{mainth}{4}

\setcounter{mainpr}{4}

\setcounter{mainex}{4}
\begin{center}
    \huge{Классическая вероятностная схема}
\end{center}
$(\Omega,\, \mc{F},\, P)\ -$ вероятностное пространство

$\Omega = \{\omega_1,\,\dots,\, \omega_n\}$, но когда в $\Omega$ конечное число элементов, как здесь, достаточно $(\Omega, P)$
\\

$P:\, \omega_i \to p_i$

$P(A) = \dfrac{\# A}{\# \Omega}\ -$ классическая формула определения вероятности
\\

$\implies p_i = \dfrac{1}{n},\, n = \# \Omega$ (все исходы равновероятны)
\begin{center}
    \huge{Независимость событий}
\end{center}
\defin{
    Пусть $(\Omega, \mc{F}, P)\ -$ вероятностное пространство. События $A_1,\dots, A_n \in \mc{F}$ называются независимыми, если для всех $1 \leq i_1 \leq \cdots \leq i_k \leq n$
    \[P\left(A_{i_1} \cap A_{i_2} \cap \cdots \cap A_{i_k}\right) = P\left(A_{i_1}\right) \cdot P\left(A_{i_2}\right) \cdots P\left(A_{i_k}\right)\]
}
\example{
    $A_1,\, A_2,\, A_3$ независимы
    \[P(A_1 \cap A_2) = P(A_1) \cdot P(A_2)\]
    \[P(A_2 \cap A_3) = P(A_2) \cdot P(A_3)\]
    \[P(A_1 \cap A_3) = P(A_1) \cdot P(A_3)\]
    \[P(A_1 \cap A_2 \cap A_3) = P(A_1) \cdot P(A_2) \cdot P(A_3)\]
    Если события попарно независимы (то есть верно $P(A_i \cap A_j) = P(A_i) \cap P(A_j)\ \forall i,j$), то из этого не следует независимость в совокупности (когда для пересечения $\geq 3\ A_i$ верно)
}
Замечание: отрицание $\left(A_i^c\right)$ не влияет на независимость
\defin{
    Пусть $(\Omega, \mc{F}, P),\ A_j, j \in J\ -$ произвольная совокупность событий. События $A_j$ называются 
    
    независимыми, если $\forall n \in \mb{N}$ любые $n$ из этих событий независимы
}
\defin{
    Условной вероятностью $P(A\mid B)$ события $A$ при условии $B$ называется
    \[P(A\mid B) = \frac{P(A \cap B)}{P(B)} \implies P(A \cap B) = P(A \mid B) \cdot P(B)\]
}
\example{
    Игральную кость кидают один раз. Какова вероятность, что выпало 2 при условии, что выпало чётное число
    \[P(''2''\mid \text{чётное}) = \frac{\dfrac{1}{6}}{\dfrac{1}{3}} = \frac{1}{2}\]
}
Замечание 1: $P(A \cap B) = P(A\mid B) \cdot P(B)\ -$ теорема умножения

Замечание 2: для независимых событий $A$ и $B$
\[P(A\mid B) = P(A)\]
\[P(B\mid A) = P(B)\]
\example{
    А $-$ авария

    П $-$ водитель пьян

    Т $-$ водитель трезв

    $P(\text{П}\mid\text{А}) = 0.4,\ P(\text{П}) = 0.05,\ P(\text{T}\mid\text{А}) = 0.6$
    Нужно решить, каким ездить, пьяным или трезвым, то есть
    \[\frac{P(\text{A}\mid \text{П})}{P(\text{А} \mid \text{Т})} = ?\]
    \\
    \[P(\text{A}\mid \text{П}) = \frac{P(\text{A} \cap \text{П})}{P(\text{П})} = \frac{P(\text{П}\mid\text{А}) \cdot P(\text{А})}{P(\text{П})}\]
    \[P(\text{А}\mid \text{Т}) = \frac{P(\text{T}\mid\text{А}) \cdot P(\text{А})}{P(\text{Т})}\]
    \[\implies \frac{P(\text{A}\mid \text{П})}{P(\text{А} \mid \text{Т})} = \frac{P(\text{П}\mid\text{А})\cdot P(\text{Т})}{P(\text{T}\mid\text{А}) \cdot P(\text{П})} = \frac{0.4 \cdot 0.95}{0.6 \cdot 0.05} \simeq 12\]
}
\end{document}